% !PRINTED 181
\weberpage{181}

\weq{9}{\varphi'(x_i) = 0 \qquad i = 1, 2, \ldots m,}

and if this determinant vanishes, then there is (\S.\,24, III) a system of quantities

\weq{10}{x_1^{(0)}, x_2^{(0)} \ldots x_m^{(0)},}

whose elements do not all vanish, and which, substituted for the $x$, satisfy the equations (9).

But now, if $x_i$, $x_i'$ are any two systems of quantities, and $\lambda$ likewise denotes an arbitrary quantity, one has the identical equations

\weq{11}{\begin{aligned}
\varphi(x + \lambda x') &= \varphi(x) + \lambda \overset{i}{\Sigma} x_i \varphi'(x_i') + \lambda^2 \varphi(x')\\
\varphi(x') &= \tfrac{1}{2}\Sigma x_i' \varphi'(x_i'),
\end{aligned}}

[\S.\,16, (12) to (15)], so that, if for $x_i'$ the values $x_i^{(0)}$ are substituted,

\weq{12}{\varphi(x_i + \lambda x_i^{(0)}) = \varphi(x),}

becomes entirely independent of $\lambda$. If we now assume that $x_1^{(0)}$ is different from zero, and we set

\[
\lambda = -\frac{x_1}{x_1^{(0)}},
\]

and further:

\weq{13}{\begin{aligned}
y_2 &= x_2 x_1^{(0)} - x_1 x_2^{(0)}\\
y_3 &= x_3 x_1^{(0)} - x_1 x_3^{(0)}\\
&\cdots\cdots\cdots\cdots\\
y_m &= x_m x_1^{(0)} - x_1 x_m^{(0)},
\end{aligned}}

then it follows that

\weq{14}{x_1^{(0)\,2} \varphi(x) = \varphi(0, y_2, y_3 \ldots y_m),}

whereby $\varphi(x)$ is expressed in terms of the $m - 1$ variables $y_2, y_3 \ldots y_m$.

\begin{center}
\S.\,57.

\gesperrt{Transformation of Quadratic Forms into a\\
Sum of Squares.}
\end{center}

We at once make an important application of the theorem just proved to the discriminant of the function (11) (regarded as a function of the second degree in $\lambda$), in which
%CONT>>
